\subsection{Generative Models}

Generative modeling learns distributions, simulators, or conditional samplers rather than only point predictions.

\begin{definition} \label{def:sl-latent-variable-generative-model} \hlt{Latent-Variable Generative Model}\\
A latent-variable model samples $Z\sim p(z)$ and then $X\sim p_{\theta}(x\mid z)$.  Learning fits $\theta$ so the marginal $p_{\theta}(x)=\int p_{\theta}(x\mid z)p(z)\,dz$ matches the data distribution.
\end{definition}

\begin{definition} \label{def:sl-variational-autoencoder} \hlt{Variational Autoencoder}\\
A variational autoencoder combines a decoder $p_{\theta}(x\mid z)$ with an encoder $q_{\phi}(z\mid x)$ and maximizes an evidence lower bound consisting of reconstruction and latent regularization terms.
\end{definition}

\begin{definition} \label{def:sl-generative-adversarial-network} \hlt{Generative Adversarial Network}\\
A generative adversarial network trains a generator against a discriminator that tries to distinguish generated samples from data.  The game defines a learned discrepancy between distributions.
\end{definition}

\begin{definition} \label{def:sl-diffusion-model} \hlt{Diffusion Model}\\
A diffusion model corrupts data through a noising process and learns a reverse denoising process.  Sampling starts from noise and iteratively applies learned denoising transitions.
\end{definition}
